\subsection{3. Exakt geschlossener Ablauf und messbarer
Unterschied}\label{exakt-geschlossener-ablauf-und-messbarer-unterschied}

Auf N=2 hat die gemeinsame Paarabbildung C\_c exakt den Gramoperator

C\_c C\_cᵀ = 8 {[}{[}I60,c²I60{]},{[}c²I60,I60{]}{]}.

Der Prüfer baut dafür die ursprünglichen 480 W-Verticesterme mit ihren
korrekten CAR-Vorzeichen auf allen C(128,2)=8128
Zweifermion-Basiszuständen aus. Insbesondere bekommt d\_i†a\_j† beim
Umsortieren das Minuszeichen. Beide lokalen Diagonalblöcke bleiben 8I60;
der neue Überlappungsblock ist 8c²I60.

Für jedes A ist der Raum aus drei normierten Paar-Flavorkombinationen
(aa, symmetrisch-ad, dd) und beiden Bosonen ein \textbf{exakter
reduzierender 5-Zustandsraum}, keine Projektion mit weggelassener
Leakage. Seine Kopplungsmatrix ist

B = √8 g {[}{[}1,0,0{]},{[}c²,√2cs,s²{]}{]}, H₅ =
{[}{[}0₃,Bᵀ{]},{[}B,ΔI₂{]}{]}.

Ein Paarzustand ist dunkel. Vom linken Boson aus ist der zyklisch
erreichbare Raum für g≠0 \textbf{genau vierdimensional}: Der
Gramdeterminant der ersten vier Krylovvektoren ist 8⁶
c⁸(1−c⁴)g¹²\textgreater0. Der gesamte N=2-Sektor hat Dimension 8248,
Paarabbildungsrang 120 und 8008 fermionische Dunkelzustände. Die 60
reduzierenden Fünferblöcke enthalten bereits 60 dieser Dunkelzustände;
die übrigen 7948 liegen außerhalb ihrer Summe.

Die beiden hellen Kopplungsquadrate sind λ±=8(1±c²). Für Ω±=√(Δ²+4g²λ±)
lautet die exakte Bosonamplitude jeder Rabi-Dublette

\codeword{u±(t)=e\textasciicircum{}(−iΔt/2)[cos(Ω±t/2)−i(Δ/Ω±)sin(Ω±t/2)].}

Damit gilt, ohne Schrieffer-Wolff- oder Klein-g-Näherung:

P(L→R;t)=\textbar(u+(t)−u−(t))/2\textbar² = 16g⁴c⁴t⁴ + O(t⁶).

Der volle G-invariante Test benötigt keine Auswahl eines intern
ausgezeichneten A: Präpariere die gleichgewichtete Mischung
ρL=(1/60)Σ\_A \textbar b\_A,L⟩⟨b\_A,L\textbar{} und lies die gesamte
rechte Bosonzahl Nb,R aus. Beide sind G-invariant. Da alle A denselben
Kanalblock haben, ergibt sich dieselbe obige Wahrscheinlichkeit. Die
Präparierbarkeit und Messbarkeit werden als Gegenmodell-Dictionary
vorausgesetzt, nicht hergeleitet.

{\def\LTcaptype{none} % do not increment counter
\begin{longtable}[]{@{}lrr@{}}
\toprule\noalign{}
Gemeinsame-Fermion-Komposition & c=3/5, s=4/5 & c=4/5, s=3/5 \\
\midrule\noalign{}
\endhead
\bottomrule\noalign{}
\endlastfoot
Lokales WWᵀ in jedem Chart & 8I60 & 8I60 \\
λ+, λ− & 272/25, 128/25 & 328/25, 72/25 \\
Koeffizient von g⁴t⁴ im Transfer & 1296/625 & 4096/625 \\
Dimensionsloses viertes Moment unten & 706/625 & 881/625 \\
\end{longtable}
}

Für dieselbe ρL seien μk=Tr(ρL H\_c\^{}k). Beide Modelle haben exakt

μ1=Δ, μ2=Δ²+8g², μ3=Δ³+16Δg².

Erst μ4=Δ⁴+24Δ²g²+64(1+c⁴)g⁴ trennt sie. Besonders wichtig:

I4 = (μ4−3μ1²μ2+2μ1⁴)/(μ2−μ1²)² = 1+c⁴,

c=(I4−1)\^{}(1/4) für die erklärte nichtnegative Familie.

\textbf{Positive Umkehrung:} Eine höhere Antwort macht den fehlenden
globalen Überlappungsparameter beobachtbar. Er kann nicht durch eine
geänderte Zeitskala oder bloße Umbenennung bei erhaltener
Vorbereitungs-/Messdictionary verschwinden. Unser Gegenmodell behauptet
ausdrücklich \textbf{nicht} identische globale N=2-Daten: Sein globaler
Gramblock unterscheidet die Komposition bereits. Der parallele
Ein-Bosonbank-Strang der Hauptrunde untersucht die stärkere Grenze, dass
auch identische vollständige N=2-Daten noch unterschiedliche
N=3-Antworten erlauben; dessen Beweis und allgemeine komplexe
Konventionen sind nicht Teil dieses unabhängigen Zertifikats.

\subsection{4. Echter gemeinsamer Ressourcenanteil, keine versteckte
Einzelkopienzerlegung}\label{echter-gemeinsamer-ressourcenanteil-keine-versteckte-einzelkopienzerlegung}

Die beiden reellen symmetrischen Flavor-Kopplungstensoren lauten CL=uuᵀ
und CR=vvᵀ, u=(1,0), v=(c,s). Sie erfüllen

Tr(CL CR)=c²∈(0,1), ‖{[}CL,CR{]}‖²\_HS=2c²s²\textgreater0.

Zwei Rang-eins-symmetrische Matrizen können hier nicht gleichzeitig
durch denselben passiven unitären Flavorbasiswechsel Takagi-diagonal
werden: Dafür müssten ihre erzeugenden Geraden parallel oder orthogonal
sein; ihr Skalarprodukt bleibt unter einem gemeinsamen Unitärwechsel
erhalten. Das ist nicht der Fall. Der gemeinsame Projektorkommutant in
M₂(C) ist exakt eindimensional, Span(I₂).

Somit gibt es \textbf{keine nichttriviale reine Einzel-Flavor-Parität
nach einem gemeinsamen passiven Basiswechsel, die beide benannten
Bosonbanken unverändert lässt}. Eine solche Parität müsste u und v mit
Vorzeichen ±1 erhalten; wegen ihrer nichtverschwindenden Überlappung
haben beide dasselbe Vorzeichen. Es bleibt nur die gesamte
Fermionparität. Die Aussage verbietet nicht zusätzliche
Transformationen, die gleichzeitig Bosonbankphasen oder Banknamen
ändern.

Das liefert den präzisen Unterschied zur Ein-Bosonbank-Familie mit einem
einzigen symmetrischen Site-Tensor: Mehrere nichtkompatible
Kopplungstensoren können eine gemeinsame Einzelkopien-Paritätszerlegung
aufbrechen. Deren \textbf{Auswahl} bleibt aber neuer Geometrieinput.

\subsection{5. Liefern frühe Quellprojektoren bereits
c?}\label{liefern-fruxfche-quellprojektoren-bereits-c}

Der neu ausgeführte Original-P0/P1-Weg bestätigt exakt

\textbar⟨zα/2,zβ/2⟩\textbar² ∈ \{0,1/4,1/2,1\}.

Die Zahlen c=1/2 oder c=1/√2 sind also als Beträge konkreter
ursprünglicher C⁴-Strahlenüberlappungen vorhanden. Dies ist mehr als
eine frei erfundene Kandidatenzahl. Es ist aber noch \textbf{kein
typkorrekter Quellschritt zu diesem H\_c}:

\begin{enumerate}
\def\labelenumi{\arabic{enumi}.}
\tightlist
\item
  Die C⁴-Projektoren wirken auf dem bestehenden Compilerträger. Das
  Flavor-C² unseres Gegenmodells ist dagegen eine G-triviale
  Multiplizitätsebene des 64-dimensionalen nativen Fermionträgers.
\item
  Die untersuchten Quellen liefern keine
  Intertwiner-/Ressourcenzuordnung, die gerade zwei dieser Strahlen zu
  den beiden globalen CAR-Chart-Einbettungen macht.
\item
  Sie wählen keine zwei Strahlen, deren Phasenlift, Bosonbankzuordnung
  oder Summe von Vertices als globalen Prozess aus.
\end{enumerate}

Das bloße Einsetzen einer bekannten Quellzahl würde diese drei Schritte
nicht ersetzen. Unsere rationalen c=3/5,4/5 sind bewusst klar benannte
Gegenmodellparameter; sie werden nicht als ursprüngliche
Strahlenüberlappungen verkauft. \textbf{Falls} künftig der fehlende
typisierte Strahlen→Chart-Schritt abgeleitet wird, kann die Familie auf
diskrete Kandidaten eingeschränkt werden und I4 deren unterschiedliche
Antworten direkt testen.

Der Typengpass lässt sich noch genauer fassen: Soweit das C⁴ in der
bereits verwendeten nativen Identifikation den internen
SU(4)-Fundamentalfaktor trägt, darf es nicht ohne weiteren Schritt als
unabhängiger \textbf{G-trivialer} Multiplizitätsraum ausgegeben werden.
Die 15 originalen SU(4)-Generatoren haben auf C⁴ keinen gemeinsamen
Fixvektor, und ihr Endomorphismenkommutant besteht exakt nur aus
Skalaren; beides ist separat geprüft. Also Hom\_G(1,4)=0 und
End(4)=1⊕\allowbreak{}15; der skalare Anteil jedes normierten Rang-eins-Projektors
ist derselbe I4/4. Ein fixer linearer equivarianter Import erzeugt
hieraus keine verschiedenen Kopien. \textbf{Paarinvarianten} wie
Tr(PαPβ) sind dagegen durchaus unterscheidbar und bleiben brauchbare
Kandidatendaten; ihnen fehlt weiterhin die erklärte globale
Einbettungs-/Vertexregel. Eine andere Feldzuordnung oder gezielte
Symmetriebrechung ist damit nicht allgemein verboten.

\subsection{6. Nächster erster fehlender
Quellschritt}\label{nuxe4chster-erster-fehlender-quellschritt}

Gesucht ist eine explizite, typisierte globale Zuordnung

Originalwort / markiertes Quellobjekt → gemeinsamer Einteilchenraum
Vglobal, Isometrien Ex:C64→Vglobal und Bosonbank-/Vertex-Zuordnung.

Sie muss insbesondere Ex†Ey festlegen. Erst dann darf ein Parameter wie
c als hergeleitet gelten. Danach sind H, erhaltene Größen und die im
Elternstrang ergänzte höhere geladene Antwort auf derselben Grundlage zu
prüfen. Die bereits vorhandenen lokalen Zahlen, vollständige interne
Symmetrie oder wiederholter Clock wählen diesen Schritt nicht von
selbst.

\subsection{7. Reproduktion und
Prüfgrenze}\label{reproduktion-und-pruxfcfgrenze}

\codeword{python3 verify.py --output results\_normal.json} und
\codeword{python3 -OO verify.py --output results\_optimized.json};
\url{REPLAY.json} versiegelt Exitcodes, Prüferhash und identische
Resultatbytes. \codeword{replay.py} führt beide Varianten lokal aus.
Keine Netz-/Live-Repo-Imports, keine große Grundzustandsdiagonalisierung
und keine Änderung älterer v1.6.8-Dateien.

Alle neuen algebraischen Aussagen sind mit ganzen Zahlen,
Gauß-Ganzzahlen oder SymPy-Rationalen geprüft; es gibt \textbf{keine
numerischen Toleranzprüfungen}. Der native Konstruktor verwendet
komplexe Maschinenarrays für ganzzahlige Gaußwerte und kontrollierte
kleine exakte Operationen; W- und globale Paargramme sind int64.
\textbf{270 eigene exakte Guards, 0 numerische Guards}, daneben 17
übernommene native Konstruktionsguards und 1073 übernommene
P0/P1-Adapterguards. Die Clock-Originalguards sind einzeln in der
eigenen Checkliste gekennzeichnet; keine Quelle wird als neuer
unabhängiger Beweis ausgegeben.

Ein erster Prüflauf wurde korrekt abgebrochen, weil SymPy den Ausdruck
2(−1−i)(−1+i) vor Expansion strukturell nicht mit 4 gleichsetzt. Ein
Minimalbeispiel reproduzierte genau diese Darstellungsursache; die
Korrektur besteht nur in exakter polynomialer Expansion vor dem
Normvergleich. Kein Toleranzwert und keine abgeschwächte mathematische
Bedingung wurden eingeführt.

\begin{center}\rule{0.5\linewidth}{0.5pt}\end{center}

\section{Teilbericht B: vereinfachter Eingriff, Aufzeichnung und
Arbeit}\label{teilbericht-b-vereinfachter-eingriff-aufzeichnung-und-arbeit}

\section{v1.6.9: Bosonzahlkontrolle ersetzt den zusätzlichen
Einzelmodenimpuls}\label{v1.6.9-bosonzahlkontrolle-ersetzt-den-zusuxe4tzlichen-einzelmodenimpuls}

\subsection{Ergebnis und genaue
Verbesserung}\label{ergebnis-und-genaue-verbesserung}

Der kausale Dreizustandsversuch benötigt keinen separat gewährten Impuls
auf Fermionmode 4. Die globale Bosonzahlphase

\[\adjustbox{max width=\linewidth}{$\displaystyle 
Z_b=e^{i\pi N_b}
$}\]

aus dem bereits dokumentierten Kontrolltyp Nb genügt für einen exakten,
nichtnachselektierten Eingriffseffekt auf die spätere Besetzung n5. Der
Tensor, das freie H, der N=2-Startzustand und der Dreizustandsraum
bleiben dieselben wie in v1.6.8.

Die Verbesserung ist algebraisch streng: Der alte Impuls Z4 liegt
außerhalb der auf diesem Raum von X und Nb erzeugten Algebra. Der neue
Impuls liegt innerhalb. Ein zusätzlicher Modenselektor wird damit
eingespart. \textbf{Das beweist noch keine vom Compiler bereitgestellte
Schaltbarkeit.} Der Unterschied zwischen fixem H und unabhängig
steuerbarem X/Nb bleibt ausdrücklich bestehen.

Ein ergänzendes Instrument zeichnet nur Paar- gegenüber Bosonzweig auf.
Es erhält die Kohärenz zwischen allen Fermionenpaaren und den
vollständigen gemeinsamen Zustand einschließlich Zeiger. Beim Weglassen
des Zeigers ändert sich die Systemkohärenz zwischen Paar und Boson; die
daraus folgende Empfängerwirkung ist exakt die Hälfte des
Phasenimpulseffekts. Die Konstruktion enthält einen wirklichen
Zeigerkoppler und ist keine Umbenennung von Messergebnissen.

\subsection{1. Was ursprünglich dokumentiert
ist}\label{was-urspruxfcnglich-dokumentiert-ist}

Acht Eingaben sind unverändert unter \codeword{inputs/} eingefroren.
\codeword{inputs\_manifest.json} enthält Pfade, Größen und SHA-256. W
bleibt
\codeword{3f00a0892157a691e4ef26920c3702772c7bf76a250fad98ff04a079bb64b763}.

Die eingefrorene \codeword{operations\_commutant.py}, Zeilen 6-18,
definiert das feste Modell H=ΔNb+gX und getrennt den Kontrolltier
A=\{X,Nb\}. Die Datei prüft in Zeilen 389-414 außerdem die Kommutation
sämtlicher 60 innerer Generatoren und der Clock mit X und Nb. Die
entsprechenden Focklifts in \codeword{native\_common.py}, Zeilen
104-219, erhalten die Besetzungszahlen. Der frühere Quellenbericht
\codeword{v168\_mechanism.md} trennt diese gewährten Kontrollstufen
bereits von einer Herleitung aktiver Geräte aus dem Compiler. Die großen
Kommutantenprogramme wurden hier nicht erneut ausgeführt.

Auf dem exakt selben Raum

\[\adjustbox{max width=\linewidth}{$\displaystyle 
\mathcal K_3=\operatorname{span}\{p_0,R_7,b_0\},\qquad p_0=|4,57\rangle,
\quad R_7=\frac1{\sqrt7}\sum_{q=1}^7s_q|p_q\rangle
$}\]

lauten die Operatoren

\[\adjustbox{max width=\linewidth}{$\displaystyle 
X_3=\begin{pmatrix}0&0&-1\\0&0&\sqrt7\\-1&\sqrt7&0\end{pmatrix},
\quad B=N_b|_{\mathcal K_3}=\operatorname{diag}(0,0,1),
\quad H_3=\Delta B+gX_3,
$}\] \[\adjustbox{max width=\linewidth}{$\displaystyle 
Z_b|_{\mathcal K_3}=I-2B=\operatorname{diag}(1,1,-1),\qquad
E=J^\dagger n_5J=\operatorname{diag}(0,1/7,0).
$}\]

Die Formeln I−2Nb und die zwei Besetzungsprojektoren unten gelten auf
N=2. Auf dem ganzen Boson-Fockraum ist Zb die Parität exp(iπNb);
außerhalb N=2 ist I−2Nb nicht dieselbe Operation.

Der dunkle Projektor ist \(P_D=I-X_3^2/8\). Die Algebra von X3 und B ist
exakt \(\mathbb C P_D\oplus M_2\), komplexe Dimension fünf. Eine Basis
ist I, B, X3, i{[}B,X3{]}, X3²; ihre Multiplikationsschließung wurde
geprüft. Alle ihre Elemente kommutieren mit PD. Für den alten
\(Z_4=\operatorname{diag}(-1,1,1)\) gilt dagegen {[}Z4,PD{]}≠0. Es wird
folglich tatsächlich ein bisher unabhängiger Operatortyp entfernt. Die
neue Kontrollklasse erzeugt aber nicht die volle M3-Algebra.

\subsection{2. Exakter unbedingter
Eingriff}\label{exakter-unbedingter-eingriff}

Setze

\[\adjustbox{max width=\linewidth}{$\displaystyle 
\Omega=\sqrt{\Delta^2/4+8g^2},\qquad
\tau=\frac\pi{2\Omega},\qquad d=\frac\Delta{2\Omega},\qquad
U=e^{-i\tau H_3}.
$}\]

Beide Arme starten im gleichen Produktzustand p0, alle Bosonen leer:

\[\adjustbox{max width=\linewidth}{$\displaystyle 
\text{frei: }U\,I\,U|p_0\rangle,
\qquad\text{Eingriff: }U\,Z_b\,U|p_0\rangle.
$}\]

Anschließend wird n5 unbedingt ausgelesen. Zb wirkt auf die bosonische
Algebra; n5 auf eine Fermionmode. Sie kommutieren auf der ganzen
ursprünglichen Fock-Algebra. Unmittelbar beim Impuls bleibt die
Empfängerstatistik für jeden Zustand unverändert. Die Entwicklung nach
dem Impuls erzeugt

\[\adjustbox{max width=\linewidth}{$\displaystyle 
p_{\rm frei}(n_5=1)=\frac{1+\cos(\pi d)}{32},
$}\] \[\adjustbox{max width=\linewidth}{$\displaystyle 
p_{Z_b}(n_5=1)=
\frac{1+(1-2d^2)^2-2(1-2d^2)\cos(\pi d)}{64},
$}\] \[\adjustbox{max width=\linewidth}{$\displaystyle 
\boxed{\delta_{b\to5}=
-\frac{(1-d^2)\left[d^2+\cos(\pi d)\right]}{16}}.
$}\]

Dies ist ein kausaler Vergleich von Eingriff und Nicht-Eingriff auf
derselben Quelle. Keine Ergebnisbedingung und keine Zweipunktkorrelation
ersetzt den Vergleich. Ein Zb direkt auf dem anfänglichen Zustand p0
hätte keinen Effekt; die native Vorentwicklung ist deshalb ein
wesentlicher Teil des Versuchs.

\subsubsection{Strikter endlicher
Nachweis}\label{strikter-endlicher-nachweis}

Für Δ\textgreater0 und 0\textless g²≤5Δ²/128 gilt 2/3≤d\textless1. Aus
\(\cos x\ge1-x^2/2\) und \(\pi<22/7\) folgt

\[\adjustbox{max width=\linewidth}{$\displaystyle 
d^2+\cos\pi d
\le -(1-d)\left[1+d-\frac{\pi^2}{2}(1-d)\right]<0.
$}\]

Denn die eckige Klammer ist größer als (291d−193)/49 und damit
mindestens 1/49. Also ist δb→5 in diesem ganzen Parameterintervall
strikt positiv. Außerhalb gilt die exakte Formel weiterhin, aber nicht
dieselbe Vorzeichenaussage.

Am Prüfpunkt Δ=1, g=1/20 sind τ≈3.02299894039036 und d=5/(3√3):

{\def\LTcaptype{none} % do not increment counter
\begin{longtable}[]{@{}lr@{}}
\toprule\noalign{}
Größe & Wert \\
\midrule\noalign{}
\endhead
\bottomrule\noalign{}
\endlastfoot
n5 ohne Impuls & 0.0002194998817122665 \\
n5 mit Bosonphase & 0.0005299169088385700 \\
Unbedingte Differenz & +0.0003104170271263035 \\
Endwert Nb ohne Impuls & 0 exakt \\
Endwert Nb mit Impuls & 25/729 exakt \\
\end{longtable}
}

Im Gegensatz zum alten Z4-Protokoll haben die beiden Endzustände also
verschiedene Zusammensetzungen. Der neue Test beweist eine kausal
beeinflusste Fermionbesetzung; er beansprucht keine in beiden Armen
bosonfreie Paarumverteilung am Endpunkt. Nb ist ein globaler bosonischer
Kontrolltyp, kein aus dem Compiler abgeleiteter räumlich lokaler Sender.

